% ======================================================================
%  AT-MONO107 — Chapter 8: The Lock Law
%  Canonical Monograph (v2.0) · The Actualization Theory:
%  A Reconstruction of Physics from Difference, Actualization and Spectrum
%
%  Status: COMPLETE — PASS-READY (MONO007 referee-ready architecture)
%  Inputs: Chapter01_Difference.tex ... Chapter07_OperatorBasis.tex
%  Method: assembly only — canonical results only, no new physics, no speculation
%  Citations: QG313 (lock universality), QG314 (organization predictor),
%             QG315 (early lock prediction), QG316 (organization phase transition),
%             QG318 (reorganization prediction), QG318.1 (lock origin),
%             QG319 (false positive audit), QG319.1 (competing predictor)
%
%  Usage: % ======================================================================
%  AT-MONO107 — Chapter 8: The Lock Law
%  Canonical Monograph (v2.0) · The Actualization Theory:
%  A Reconstruction of Physics from Difference, Actualization and Spectrum
%
%  Status: COMPLETE — PASS-READY (MONO007 referee-ready architecture)
%  Inputs: Chapter01_Difference.tex ... Chapter07_OperatorBasis.tex
%  Method: assembly only — canonical results only, no new physics, no speculation
%  Citations: QG313 (lock universality), QG314 (organization predictor),
%             QG315 (early lock prediction), QG316 (organization phase transition),
%             QG318 (reorganization prediction), QG318.1 (lock origin),
%             QG319 (false positive audit), QG319.1 (competing predictor)
%
%  Usage: % ======================================================================
%  AT-MONO107 — Chapter 8: The Lock Law
%  Canonical Monograph (v2.0) · The Actualization Theory:
%  A Reconstruction of Physics from Difference, Actualization and Spectrum
%
%  Status: COMPLETE — PASS-READY (MONO007 referee-ready architecture)
%  Inputs: Chapter01_Difference.tex ... Chapter07_OperatorBasis.tex
%  Method: assembly only — canonical results only, no new physics, no speculation
%  Citations: QG313 (lock universality), QG314 (organization predictor),
%             QG315 (early lock prediction), QG316 (organization phase transition),
%             QG318 (reorganization prediction), QG318.1 (lock origin),
%             QG319 (false positive audit), QG319.1 (competing predictor)
%
%  Usage: % ======================================================================
%  AT-MONO107 — Chapter 8: The Lock Law
%  Canonical Monograph (v2.0) · The Actualization Theory:
%  A Reconstruction of Physics from Difference, Actualization and Spectrum
%
%  Status: COMPLETE — PASS-READY (MONO007 referee-ready architecture)
%  Inputs: Chapter01_Difference.tex ... Chapter07_OperatorBasis.tex
%  Method: assembly only — canonical results only, no new physics, no speculation
%  Citations: QG313 (lock universality), QG314 (organization predictor),
%             QG315 (early lock prediction), QG316 (organization phase transition),
%             QG318 (reorganization prediction), QG318.1 (lock origin),
%             QG319 (false positive audit), QG319.1 (competing predictor)
%
%  Usage: \input{Chapter08_LockLaw.tex} after the preamble
%         that defines the theorem environments (or include via the master file).
% ======================================================================

% ---- Theorem environments (define in the master preamble if not present) ----
% \newtheorem{definition}{Definition}[chapter]
% \newtheorem{lemma}[definition]{Lemma}
% \newtheorem{theorem}[definition]{Theorem}
% \newtheorem{corollary}[definition]{Corollary}
% \newtheorem{remark}[definition]{Remark}

\chapter{The Lock Law}
\label{c08:chap:locks}

\section{Introduction}
\label{c08:sec:locks-introduction}

In Chapter~\ref{c07:chap:operators} we established the operator basis of the theory: the four
operators Crowding, Compression, Beat, and Locking, derived as projections of the D96
spectrum. This chapter develops the deeper structure of the Locking operator --- the
\emph{lock law}: the lock identities of organized spectra carry a distinctive structure ---
the lock \emph{law} is universal (the presence of stable, reproducible characteristic
ratios), while the lock \emph{values} are domain-specific \cite{c08:QG313}. We establish the
temporal and structural properties of the locks --- that they precede organization
\cite{c08:QG315}, that organization emerges at a critical transition \cite{c08:QG316}, and
that the lock structure traces to the moment-chain identity \cite{c08:QG318.1}.

The lock results are established on \emph{synthetic deterministic evolving-law cohorts}
\cite{c08:QG315,c08:QG316}: the claims are cohort-scoped and the synthetic basis is disclosed,
per the referee-safe requirements of MONO007 \cite{c08:MONO007}. The honest scope of the lock
rule is reported: competing standard complexity measures match or beat the lock rule on the
evolving cohort \cite{c08:QG319.1}, and the lock rule is weak as a standalone detector under
adversarial synthesis \cite{c08:QG319}. These results scope, rather than contradict, the
lock law.

\section{Locking as an Operator}
\label{c08:sec:locking-operator}

\begin{theorem}[Locks are Locking projections]
\label{c08:thm:locks-locking}
A \emph{lock identity} is a normalized ratio of a spectrum's moments --- the moment/span,
compression/count, higher-moment, and $\sqrt{\text{moment}}/\text{span}$ ratios --- that
takes a stable, reproducible value characteristic of the spectrum's organization. The lock
identities are projections of the Locking operator: normalized ratio reads of the spectrum
\cite{c08:QG313}, consistent with the operator basis of Chapter~\ref{c07:chap:operators}
(Theorem~\ref{c07:thm:operators-derived}). They are derived reads, not primitives: by
Chapter~\ref{c07:chap:operators} (Corollary~\ref{c07:cor:operators-not-primitive}) the
operators are derived projections of the spectrum, and by Chapter~\ref{c06:chap:d96}
(Theorem~\ref{c06:thm:d96-derived}) the spectrum is the derived output of the attractor ---
the locks are thrice-removed from primitiveness.
\end{theorem}

\section{The Universal Lock Law}
\label{c08:sec:universal-law}

\begin{theorem}[The lock law is universal in structure]
\label{c08:thm:lock-law-universal}
The lock law is universal: every organized domain carries stable, reproducible lock
identities. The \emph{presence} of lock structure recurs across all organized domains.
\end{theorem}

\begin{proof}
The lock-universality audit \cite{c08:QG313} measures the normalized lock identities across
seven domains --- physics (D96 multiplicities), language (Zipf), music (harmonic
octave-class), DNA (codon inequality), software (token power law), finance (heavy tail),
networks (connectome) --- and finds stable lock structure in all seven: the lock law recurs
in every organized domain. Locks are robust content: unlike the binary operator presence,
which can be faked --- the operator basis can be partially faked by two-level crafted
spectra \cite{c08:QG312}, but the beat-identity locks are carried by none of the fakes ---
the lock identities are not trivially reproducible by adversarial spectra. The locks are
therefore the robust organization content, in contrast to the weak binary screen, and this
is the basis of the lock law's status as the deeper organization signature.
\end{proof}

\section{Domain-Specific Lock Values}
\label{c08:sec:domain-values}

\begin{theorem}[Lock values are domain-specific]
\label{c08:thm:lock-values-domain}
The lock \emph{values} are domain-specific: the characteristic ratios differ across domains,
each domain carrying its own lock values.
\end{theorem}

\begin{proof}
The lock-universality audit \cite{c08:QG313} measures the moment/span, compression/count,
higher-moment, and $\sqrt{\text{moment}}/\text{span}$ ratios per domain and finds distinct
values (at least four distinct moment/span values appear):
\begin{center}
\begin{tabular}{lcccc}
\hline
Domain & moment/span & comp./count & higher-moment & $\sqrt{\text{moment}}/\text{span}$ \\
\hline
physics & 31.7 & 2.41 & 2.96 & 21.4 \\
language & 45.0 & 180.6 & 369.8 & 5.70 \\
music & 19.7 & 22.1 & 26.3 & 4.67 \\
DNA & 36.0 & 5.67 & 6.35 & 16.3 \\
software & 9.19 & 523.2 & 877.7 & 0.84 \\
finance & 1.62 & 334.5 & 469.9 & 0.18 \\
networks & 66.3 & 3.97 & 4.02 & 33.5 \\
\hline
\end{tabular}
\end{center}
Hence the lock law has two faces: by Theorem~\ref{c08:thm:lock-law-universal} the structure
(the presence of stable ratios) is universal; by this theorem the values (the specific
ratios) are the domain's fingerprint --- universal in structure, specific in value.
\end{proof}

\section{Locks and Organization}
\label{c08:sec:locks-organization}

\begin{theorem}[Locks precede organization]
\label{c08:thm:locks-precede}
The lock identities precede mature organization: in evolving systems, the lock structure
appears before the full organization develops, and systems with early lock structure are
rigidity-committed: they reorganize less than un-locked systems under structural change.
\end{theorem}

\begin{proof}
The early-lock-prediction audit \cite{c08:QG315} evolves four systems (software history,
wiki edits, citation networks, language corpora) from flat to mature laws and measures lock
coherence and organization maturity at each stage: in most systems the lock identity reaches
half-strength at an earlier stage than the maturity does, and no system shows the locks
lagging the maturity. Hence the locks precede organization --- an early signature of the
developing structure. They are also a rigidity signal: the reorganization-prediction audit
\cite{c08:QG318} evolves systems through a reorganization and measures future topology
change --- systems with early lock coherence reorganize small (the locked group is 100\%
small-reorganization), while un-locked systems remain plastic and reorganize large.
\end{proof}

\section{The Critical Transition}
\label{c08:sec:critical-transition}

\begin{theorem}[Organization is a phase transition]
\label{c08:thm:phase-transition}
The organization structure emerges at a critical transition: the binary operator basis
completes discontinuously at a critical parameter, while the quantitative organization
grows continuously.
\end{theorem}

\begin{proof}
The organization-phase-transition audit \cite{c08:QG316} sweeps a continuous organization
parameter across the four regimes (white noise, weak, medium, strong), measuring the
operator basis, the lock coherence, and the organization maturity. The binary all-four
operator screen flips from incomplete to complete at the critical parameter $g^\ast \approx
0.31$ and persists for all stronger organization; the lock coherence emerges in the same
critical window; the maturity grows continuously. Hence the binary operator basis is a
critical order parameter --- a phase transition --- while the quantitative organization is
gradual, and the locks appear at the critical onset of organization, consistent with the
locks-precede result (Theorem~\ref{c08:thm:locks-precede}).
\end{proof}

\section{Predictive Consequences}
\label{c08:sec:predictive}

\begin{theorem}[The blind prediction protocol]
\label{c08:thm:blind-prediction}
The lock structure carries forward-looking information: the future organization class can be
predicted from the early lock structure, with the prediction rule fixed before the reveal.
\end{theorem}

\begin{proof}
The blind-organization-prediction result predicts the future high-maturity class from the
early-stage lock coherence, with the prediction rule fixed before the later stage is
revealed; the early lock presence identifies the future top-class systems exactly.
\end{proof}

The organization-predictor audit \cite{c08:QG314} computes an organization score from the
lock structure only and tests it against the eight domains: the score separates the
organized class from the unorganized class --- the organized systems lock above both
unorganized systems --- but does not rank the strength within the class: heavy-tailed
systems have large characteristic numerators that never lock onto a small fraction, scoring
below the Zipf systems despite being stronger organizations. Hence the lock values predict
the class, not the strength: the lock structure is a class-level predictor of organization,
distinguishing organized from unorganized systems early and reliably, without ranking them
internally.

\section{Validation Status}
\label{c08:sec:validation}

\begin{theorem}[The lock origin is the moment-chain identity]
\label{c08:thm:lock-origin}
The lock structure traces to the moment-chain identity of the D96 spectrum:
$\mathrm{occMom}/\Sigma m = (\Sigma m^2/\Sigma m)\cdot(\mathrm{occMom}/\Sigma m^2)$
holds exactly, linking the lock identities by an algebraic necessity.
\end{theorem}

\begin{proof}
The lock-origin audit \cite{c08:QG318.1} establishes the moment-chain (telescoping)
identity: the lock1 ratio $\mathrm{occMom}/\Sigma m = 20.0026$ equals the product of lock2
$\Sigma m^2/\Sigma m = 2.4105$ and lock3 $\mathrm{occMom}/\Sigma m^2 = 8.2980$, exactly.
The locks are therefore ratios of one moment hierarchy, linked by an algebraic necessity ---
the common source of lock formation: a resonance fixed point in structure, D96-specific in
value.
\end{proof}

The lock rule is cohort-scoped: the false-positive audit \cite{c08:QG319} finds the lock
rule weak as a standalone detector under adversarial synthesis --- the lock identity is
fake-able by tiny rational-ratio spectra and miss-able by large-numerator organizations ---
and the competing-predictor audit \cite{c08:QG319.1} evaluates entropy, gini, exponent,
spectral gap, and lock coherence with an identical protocol on the evolving cohort, finding
the standard measures reach the same or better accuracy. This scoping limits the lock
\emph{rule} as a detector, not the lock \emph{law} as a structural result.

\section{Summary}
\label{c08:sec:locks-summary}

This chapter has established the lock law of the theory: the lock identities are Locking
projections of the spectrum, derived and not primitive (Theorem~\ref{c08:thm:locks-locking}),
universal in structure across organized domains and robust content
(Theorem~\ref{c08:thm:lock-law-universal}), with domain-specific fingerprint values
(Theorem~\ref{c08:thm:lock-values-domain}). The locks precede organization and indicate
future rigidity (Theorem~\ref{c08:thm:locks-precede}); organization emerges at a critical
transition, the locks appearing at the critical window (Theorem~\ref{c08:thm:phase-transition});
and the lock structure predicts the future organization class, though not the strength
within it (Theorem~\ref{c08:thm:blind-prediction}). The lock origin is the moment-chain
identity of the D96 spectrum (Theorem~\ref{c08:thm:lock-origin}), and the lock rule is
cohort-scoped --- a scoping that limits the detector, not the structural law. The following
chapters read the physics sectors of the theory, each as a structured read of the D96
spectrum through the operator layer and its lock structure.

\begin{thebibliography}{99}
\bibitem{c08:QG312} AT-QG Phase 312, \emph{Adversarial Spectrum Audit}, Docs/Research.
\bibitem{c08:QG313} AT-QG Phase 313, \emph{Lock Universality Audit}, Docs/Research.
\bibitem{c08:QG314} AT-QG Phase 314, \emph{Organization Predictor Audit}, Docs/Research.
\bibitem{c08:QG315} AT-QG Phase 315, \emph{Early Lock Prediction}, Docs/Research.
\bibitem{c08:QG316} AT-QG Phase 316, \emph{Organization Phase Transition Audit}, Docs/Research.
\bibitem{c08:QG318} AT-QG Phase 318, \emph{Reorganization Prediction}, Docs/Research.
\bibitem{c08:QG318.1} AT-QG Phase 318.1, \emph{Lock Origin Audit}, Docs/Research.
\bibitem{c08:QG319} AT-QG Phase 319, \emph{False Positive Audit}, Docs/Research.
\bibitem{c08:QG319.1} AT-QG Phase 319.1, \emph{Competing Predictor Audit}, Docs/Research.
\bibitem{c08:MONO007} AT-MONO007, \emph{Referee-Readiness Resolution}, Docs/Zenodo/Monograph.
\end{thebibliography}
 after the preamble
%         that defines the theorem environments (or include via the master file).
% ======================================================================

% ---- Theorem environments (define in the master preamble if not present) ----
% \newtheorem{definition}{Definition}[chapter]
% \newtheorem{lemma}[definition]{Lemma}
% \newtheorem{theorem}[definition]{Theorem}
% \newtheorem{corollary}[definition]{Corollary}
% \newtheorem{remark}[definition]{Remark}

\chapter{The Lock Law}
\label{c08:chap:locks}

\section{Introduction}
\label{c08:sec:locks-introduction}

In Chapter~\ref{c07:chap:operators} we established the operator basis of the theory: the four
operators Crowding, Compression, Beat, and Locking, derived as projections of the D96
spectrum. This chapter develops the deeper structure of the Locking operator --- the
\emph{lock law}: the lock identities of organized spectra carry a distinctive structure ---
the lock \emph{law} is universal (the presence of stable, reproducible characteristic
ratios), while the lock \emph{values} are domain-specific \cite{c08:QG313}. We establish the
temporal and structural properties of the locks --- that they precede organization
\cite{c08:QG315}, that organization emerges at a critical transition \cite{c08:QG316}, and
that the lock structure traces to the moment-chain identity \cite{c08:QG318.1}.

The lock results are established on \emph{synthetic deterministic evolving-law cohorts}
\cite{c08:QG315,c08:QG316}: the claims are cohort-scoped and the synthetic basis is disclosed,
per the referee-safe requirements of MONO007 \cite{c08:MONO007}. The honest scope of the lock
rule is reported: competing standard complexity measures match or beat the lock rule on the
evolving cohort \cite{c08:QG319.1}, and the lock rule is weak as a standalone detector under
adversarial synthesis \cite{c08:QG319}. These results scope, rather than contradict, the
lock law.

\section{Locking as an Operator}
\label{c08:sec:locking-operator}

\begin{theorem}[Locks are Locking projections]
\label{c08:thm:locks-locking}
A \emph{lock identity} is a normalized ratio of a spectrum's moments --- the moment/span,
compression/count, higher-moment, and $\sqrt{\text{moment}}/\text{span}$ ratios --- that
takes a stable, reproducible value characteristic of the spectrum's organization. The lock
identities are projections of the Locking operator: normalized ratio reads of the spectrum
\cite{c08:QG313}, consistent with the operator basis of Chapter~\ref{c07:chap:operators}
(Theorem~\ref{c07:thm:operators-derived}). They are derived reads, not primitives: by
Chapter~\ref{c07:chap:operators} (Corollary~\ref{c07:cor:operators-not-primitive}) the
operators are derived projections of the spectrum, and by Chapter~\ref{c06:chap:d96}
(Theorem~\ref{c06:thm:d96-derived}) the spectrum is the derived output of the attractor ---
the locks are thrice-removed from primitiveness.
\end{theorem}

\section{The Universal Lock Law}
\label{c08:sec:universal-law}

\begin{theorem}[The lock law is universal in structure]
\label{c08:thm:lock-law-universal}
The lock law is universal: every organized domain carries stable, reproducible lock
identities. The \emph{presence} of lock structure recurs across all organized domains.
\end{theorem}

\begin{proof}
The lock-universality audit \cite{c08:QG313} measures the normalized lock identities across
seven domains --- physics (D96 multiplicities), language (Zipf), music (harmonic
octave-class), DNA (codon inequality), software (token power law), finance (heavy tail),
networks (connectome) --- and finds stable lock structure in all seven: the lock law recurs
in every organized domain. Locks are robust content: unlike the binary operator presence,
which can be faked --- the operator basis can be partially faked by two-level crafted
spectra \cite{c08:QG312}, but the beat-identity locks are carried by none of the fakes ---
the lock identities are not trivially reproducible by adversarial spectra. The locks are
therefore the robust organization content, in contrast to the weak binary screen, and this
is the basis of the lock law's status as the deeper organization signature.
\end{proof}

\section{Domain-Specific Lock Values}
\label{c08:sec:domain-values}

\begin{theorem}[Lock values are domain-specific]
\label{c08:thm:lock-values-domain}
The lock \emph{values} are domain-specific: the characteristic ratios differ across domains,
each domain carrying its own lock values.
\end{theorem}

\begin{proof}
The lock-universality audit \cite{c08:QG313} measures the moment/span, compression/count,
higher-moment, and $\sqrt{\text{moment}}/\text{span}$ ratios per domain and finds distinct
values (at least four distinct moment/span values appear):
\begin{center}
\begin{tabular}{lcccc}
\hline
Domain & moment/span & comp./count & higher-moment & $\sqrt{\text{moment}}/\text{span}$ \\
\hline
physics & 31.7 & 2.41 & 2.96 & 21.4 \\
language & 45.0 & 180.6 & 369.8 & 5.70 \\
music & 19.7 & 22.1 & 26.3 & 4.67 \\
DNA & 36.0 & 5.67 & 6.35 & 16.3 \\
software & 9.19 & 523.2 & 877.7 & 0.84 \\
finance & 1.62 & 334.5 & 469.9 & 0.18 \\
networks & 66.3 & 3.97 & 4.02 & 33.5 \\
\hline
\end{tabular}
\end{center}
Hence the lock law has two faces: by Theorem~\ref{c08:thm:lock-law-universal} the structure
(the presence of stable ratios) is universal; by this theorem the values (the specific
ratios) are the domain's fingerprint --- universal in structure, specific in value.
\end{proof}

\section{Locks and Organization}
\label{c08:sec:locks-organization}

\begin{theorem}[Locks precede organization]
\label{c08:thm:locks-precede}
The lock identities precede mature organization: in evolving systems, the lock structure
appears before the full organization develops, and systems with early lock structure are
rigidity-committed: they reorganize less than un-locked systems under structural change.
\end{theorem}

\begin{proof}
The early-lock-prediction audit \cite{c08:QG315} evolves four systems (software history,
wiki edits, citation networks, language corpora) from flat to mature laws and measures lock
coherence and organization maturity at each stage: in most systems the lock identity reaches
half-strength at an earlier stage than the maturity does, and no system shows the locks
lagging the maturity. Hence the locks precede organization --- an early signature of the
developing structure. They are also a rigidity signal: the reorganization-prediction audit
\cite{c08:QG318} evolves systems through a reorganization and measures future topology
change --- systems with early lock coherence reorganize small (the locked group is 100\%
small-reorganization), while un-locked systems remain plastic and reorganize large.
\end{proof}

\section{The Critical Transition}
\label{c08:sec:critical-transition}

\begin{theorem}[Organization is a phase transition]
\label{c08:thm:phase-transition}
The organization structure emerges at a critical transition: the binary operator basis
completes discontinuously at a critical parameter, while the quantitative organization
grows continuously.
\end{theorem}

\begin{proof}
The organization-phase-transition audit \cite{c08:QG316} sweeps a continuous organization
parameter across the four regimes (white noise, weak, medium, strong), measuring the
operator basis, the lock coherence, and the organization maturity. The binary all-four
operator screen flips from incomplete to complete at the critical parameter $g^\ast \approx
0.31$ and persists for all stronger organization; the lock coherence emerges in the same
critical window; the maturity grows continuously. Hence the binary operator basis is a
critical order parameter --- a phase transition --- while the quantitative organization is
gradual, and the locks appear at the critical onset of organization, consistent with the
locks-precede result (Theorem~\ref{c08:thm:locks-precede}).
\end{proof}

\section{Predictive Consequences}
\label{c08:sec:predictive}

\begin{theorem}[The blind prediction protocol]
\label{c08:thm:blind-prediction}
The lock structure carries forward-looking information: the future organization class can be
predicted from the early lock structure, with the prediction rule fixed before the reveal.
\end{theorem}

\begin{proof}
The blind-organization-prediction result predicts the future high-maturity class from the
early-stage lock coherence, with the prediction rule fixed before the later stage is
revealed; the early lock presence identifies the future top-class systems exactly.
\end{proof}

The organization-predictor audit \cite{c08:QG314} computes an organization score from the
lock structure only and tests it against the eight domains: the score separates the
organized class from the unorganized class --- the organized systems lock above both
unorganized systems --- but does not rank the strength within the class: heavy-tailed
systems have large characteristic numerators that never lock onto a small fraction, scoring
below the Zipf systems despite being stronger organizations. Hence the lock values predict
the class, not the strength: the lock structure is a class-level predictor of organization,
distinguishing organized from unorganized systems early and reliably, without ranking them
internally.

\section{Validation Status}
\label{c08:sec:validation}

\begin{theorem}[The lock origin is the moment-chain identity]
\label{c08:thm:lock-origin}
The lock structure traces to the moment-chain identity of the D96 spectrum:
$\mathrm{occMom}/\Sigma m = (\Sigma m^2/\Sigma m)\cdot(\mathrm{occMom}/\Sigma m^2)$
holds exactly, linking the lock identities by an algebraic necessity.
\end{theorem}

\begin{proof}
The lock-origin audit \cite{c08:QG318.1} establishes the moment-chain (telescoping)
identity: the lock1 ratio $\mathrm{occMom}/\Sigma m = 20.0026$ equals the product of lock2
$\Sigma m^2/\Sigma m = 2.4105$ and lock3 $\mathrm{occMom}/\Sigma m^2 = 8.2980$, exactly.
The locks are therefore ratios of one moment hierarchy, linked by an algebraic necessity ---
the common source of lock formation: a resonance fixed point in structure, D96-specific in
value.
\end{proof}

The lock rule is cohort-scoped: the false-positive audit \cite{c08:QG319} finds the lock
rule weak as a standalone detector under adversarial synthesis --- the lock identity is
fake-able by tiny rational-ratio spectra and miss-able by large-numerator organizations ---
and the competing-predictor audit \cite{c08:QG319.1} evaluates entropy, gini, exponent,
spectral gap, and lock coherence with an identical protocol on the evolving cohort, finding
the standard measures reach the same or better accuracy. This scoping limits the lock
\emph{rule} as a detector, not the lock \emph{law} as a structural result.

\section{Summary}
\label{c08:sec:locks-summary}

This chapter has established the lock law of the theory: the lock identities are Locking
projections of the spectrum, derived and not primitive (Theorem~\ref{c08:thm:locks-locking}),
universal in structure across organized domains and robust content
(Theorem~\ref{c08:thm:lock-law-universal}), with domain-specific fingerprint values
(Theorem~\ref{c08:thm:lock-values-domain}). The locks precede organization and indicate
future rigidity (Theorem~\ref{c08:thm:locks-precede}); organization emerges at a critical
transition, the locks appearing at the critical window (Theorem~\ref{c08:thm:phase-transition});
and the lock structure predicts the future organization class, though not the strength
within it (Theorem~\ref{c08:thm:blind-prediction}). The lock origin is the moment-chain
identity of the D96 spectrum (Theorem~\ref{c08:thm:lock-origin}), and the lock rule is
cohort-scoped --- a scoping that limits the detector, not the structural law. The following
chapters read the physics sectors of the theory, each as a structured read of the D96
spectrum through the operator layer and its lock structure.

\begin{thebibliography}{99}
\bibitem{c08:QG312} AT-QG Phase 312, \emph{Adversarial Spectrum Audit}, Docs/Research.
\bibitem{c08:QG313} AT-QG Phase 313, \emph{Lock Universality Audit}, Docs/Research.
\bibitem{c08:QG314} AT-QG Phase 314, \emph{Organization Predictor Audit}, Docs/Research.
\bibitem{c08:QG315} AT-QG Phase 315, \emph{Early Lock Prediction}, Docs/Research.
\bibitem{c08:QG316} AT-QG Phase 316, \emph{Organization Phase Transition Audit}, Docs/Research.
\bibitem{c08:QG318} AT-QG Phase 318, \emph{Reorganization Prediction}, Docs/Research.
\bibitem{c08:QG318.1} AT-QG Phase 318.1, \emph{Lock Origin Audit}, Docs/Research.
\bibitem{c08:QG319} AT-QG Phase 319, \emph{False Positive Audit}, Docs/Research.
\bibitem{c08:QG319.1} AT-QG Phase 319.1, \emph{Competing Predictor Audit}, Docs/Research.
\bibitem{c08:MONO007} AT-MONO007, \emph{Referee-Readiness Resolution}, Docs/Zenodo/Monograph.
\end{thebibliography}
 after the preamble
%         that defines the theorem environments (or include via the master file).
% ======================================================================

% ---- Theorem environments (define in the master preamble if not present) ----
% \newtheorem{definition}{Definition}[chapter]
% \newtheorem{lemma}[definition]{Lemma}
% \newtheorem{theorem}[definition]{Theorem}
% \newtheorem{corollary}[definition]{Corollary}
% \newtheorem{remark}[definition]{Remark}

\chapter{The Lock Law}
\label{c08:chap:locks}

\section{Introduction}
\label{c08:sec:locks-introduction}

In Chapter~\ref{c07:chap:operators} we established the operator basis of the theory: the four
operators Crowding, Compression, Beat, and Locking, derived as projections of the D96
spectrum. This chapter develops the deeper structure of the Locking operator --- the
\emph{lock law}: the lock identities of organized spectra carry a distinctive structure ---
the lock \emph{law} is universal (the presence of stable, reproducible characteristic
ratios), while the lock \emph{values} are domain-specific \cite{c08:QG313}. We establish the
temporal and structural properties of the locks --- that they precede organization
\cite{c08:QG315}, that organization emerges at a critical transition \cite{c08:QG316}, and
that the lock structure traces to the moment-chain identity \cite{c08:QG318.1}.

The lock results are established on \emph{synthetic deterministic evolving-law cohorts}
\cite{c08:QG315,c08:QG316}: the claims are cohort-scoped and the synthetic basis is disclosed,
per the referee-safe requirements of MONO007 \cite{c08:MONO007}. The honest scope of the lock
rule is reported: competing standard complexity measures match or beat the lock rule on the
evolving cohort \cite{c08:QG319.1}, and the lock rule is weak as a standalone detector under
adversarial synthesis \cite{c08:QG319}. These results scope, rather than contradict, the
lock law.

\section{Locking as an Operator}
\label{c08:sec:locking-operator}

\begin{theorem}[Locks are Locking projections]
\label{c08:thm:locks-locking}
A \emph{lock identity} is a normalized ratio of a spectrum's moments --- the moment/span,
compression/count, higher-moment, and $\sqrt{\text{moment}}/\text{span}$ ratios --- that
takes a stable, reproducible value characteristic of the spectrum's organization. The lock
identities are projections of the Locking operator: normalized ratio reads of the spectrum
\cite{c08:QG313}, consistent with the operator basis of Chapter~\ref{c07:chap:operators}
(Theorem~\ref{c07:thm:operators-derived}). They are derived reads, not primitives: by
Chapter~\ref{c07:chap:operators} (Corollary~\ref{c07:cor:operators-not-primitive}) the
operators are derived projections of the spectrum, and by Chapter~\ref{c06:chap:d96}
(Theorem~\ref{c06:thm:d96-derived}) the spectrum is the derived output of the attractor ---
the locks are thrice-removed from primitiveness.
\end{theorem}

\section{The Universal Lock Law}
\label{c08:sec:universal-law}

\begin{theorem}[The lock law is universal in structure]
\label{c08:thm:lock-law-universal}
The lock law is universal: every organized domain carries stable, reproducible lock
identities. The \emph{presence} of lock structure recurs across all organized domains.
\end{theorem}

\begin{proof}
The lock-universality audit \cite{c08:QG313} measures the normalized lock identities across
seven domains --- physics (D96 multiplicities), language (Zipf), music (harmonic
octave-class), DNA (codon inequality), software (token power law), finance (heavy tail),
networks (connectome) --- and finds stable lock structure in all seven: the lock law recurs
in every organized domain. Locks are robust content: unlike the binary operator presence,
which can be faked --- the operator basis can be partially faked by two-level crafted
spectra \cite{c08:QG312}, but the beat-identity locks are carried by none of the fakes ---
the lock identities are not trivially reproducible by adversarial spectra. The locks are
therefore the robust organization content, in contrast to the weak binary screen, and this
is the basis of the lock law's status as the deeper organization signature.
\end{proof}

\section{Domain-Specific Lock Values}
\label{c08:sec:domain-values}

\begin{theorem}[Lock values are domain-specific]
\label{c08:thm:lock-values-domain}
The lock \emph{values} are domain-specific: the characteristic ratios differ across domains,
each domain carrying its own lock values.
\end{theorem}

\begin{proof}
The lock-universality audit \cite{c08:QG313} measures the moment/span, compression/count,
higher-moment, and $\sqrt{\text{moment}}/\text{span}$ ratios per domain and finds distinct
values (at least four distinct moment/span values appear):
\begin{center}
\begin{tabular}{lcccc}
\hline
Domain & moment/span & comp./count & higher-moment & $\sqrt{\text{moment}}/\text{span}$ \\
\hline
physics & 31.7 & 2.41 & 2.96 & 21.4 \\
language & 45.0 & 180.6 & 369.8 & 5.70 \\
music & 19.7 & 22.1 & 26.3 & 4.67 \\
DNA & 36.0 & 5.67 & 6.35 & 16.3 \\
software & 9.19 & 523.2 & 877.7 & 0.84 \\
finance & 1.62 & 334.5 & 469.9 & 0.18 \\
networks & 66.3 & 3.97 & 4.02 & 33.5 \\
\hline
\end{tabular}
\end{center}
Hence the lock law has two faces: by Theorem~\ref{c08:thm:lock-law-universal} the structure
(the presence of stable ratios) is universal; by this theorem the values (the specific
ratios) are the domain's fingerprint --- universal in structure, specific in value.
\end{proof}

\section{Locks and Organization}
\label{c08:sec:locks-organization}

\begin{theorem}[Locks precede organization]
\label{c08:thm:locks-precede}
The lock identities precede mature organization: in evolving systems, the lock structure
appears before the full organization develops, and systems with early lock structure are
rigidity-committed: they reorganize less than un-locked systems under structural change.
\end{theorem}

\begin{proof}
The early-lock-prediction audit \cite{c08:QG315} evolves four systems (software history,
wiki edits, citation networks, language corpora) from flat to mature laws and measures lock
coherence and organization maturity at each stage: in most systems the lock identity reaches
half-strength at an earlier stage than the maturity does, and no system shows the locks
lagging the maturity. Hence the locks precede organization --- an early signature of the
developing structure. They are also a rigidity signal: the reorganization-prediction audit
\cite{c08:QG318} evolves systems through a reorganization and measures future topology
change --- systems with early lock coherence reorganize small (the locked group is 100\%
small-reorganization), while un-locked systems remain plastic and reorganize large.
\end{proof}

\section{The Critical Transition}
\label{c08:sec:critical-transition}

\begin{theorem}[Organization is a phase transition]
\label{c08:thm:phase-transition}
The organization structure emerges at a critical transition: the binary operator basis
completes discontinuously at a critical parameter, while the quantitative organization
grows continuously.
\end{theorem}

\begin{proof}
The organization-phase-transition audit \cite{c08:QG316} sweeps a continuous organization
parameter across the four regimes (white noise, weak, medium, strong), measuring the
operator basis, the lock coherence, and the organization maturity. The binary all-four
operator screen flips from incomplete to complete at the critical parameter $g^\ast \approx
0.31$ and persists for all stronger organization; the lock coherence emerges in the same
critical window; the maturity grows continuously. Hence the binary operator basis is a
critical order parameter --- a phase transition --- while the quantitative organization is
gradual, and the locks appear at the critical onset of organization, consistent with the
locks-precede result (Theorem~\ref{c08:thm:locks-precede}).
\end{proof}

\section{Predictive Consequences}
\label{c08:sec:predictive}

\begin{theorem}[The blind prediction protocol]
\label{c08:thm:blind-prediction}
The lock structure carries forward-looking information: the future organization class can be
predicted from the early lock structure, with the prediction rule fixed before the reveal.
\end{theorem}

\begin{proof}
The blind-organization-prediction result predicts the future high-maturity class from the
early-stage lock coherence, with the prediction rule fixed before the later stage is
revealed; the early lock presence identifies the future top-class systems exactly.
\end{proof}

The organization-predictor audit \cite{c08:QG314} computes an organization score from the
lock structure only and tests it against the eight domains: the score separates the
organized class from the unorganized class --- the organized systems lock above both
unorganized systems --- but does not rank the strength within the class: heavy-tailed
systems have large characteristic numerators that never lock onto a small fraction, scoring
below the Zipf systems despite being stronger organizations. Hence the lock values predict
the class, not the strength: the lock structure is a class-level predictor of organization,
distinguishing organized from unorganized systems early and reliably, without ranking them
internally.

\section{Validation Status}
\label{c08:sec:validation}

\begin{theorem}[The lock origin is the moment-chain identity]
\label{c08:thm:lock-origin}
The lock structure traces to the moment-chain identity of the D96 spectrum:
$\mathrm{occMom}/\Sigma m = (\Sigma m^2/\Sigma m)\cdot(\mathrm{occMom}/\Sigma m^2)$
holds exactly, linking the lock identities by an algebraic necessity.
\end{theorem}

\begin{proof}
The lock-origin audit \cite{c08:QG318.1} establishes the moment-chain (telescoping)
identity: the lock1 ratio $\mathrm{occMom}/\Sigma m = 20.0026$ equals the product of lock2
$\Sigma m^2/\Sigma m = 2.4105$ and lock3 $\mathrm{occMom}/\Sigma m^2 = 8.2980$, exactly.
The locks are therefore ratios of one moment hierarchy, linked by an algebraic necessity ---
the common source of lock formation: a resonance fixed point in structure, D96-specific in
value.
\end{proof}

The lock rule is cohort-scoped: the false-positive audit \cite{c08:QG319} finds the lock
rule weak as a standalone detector under adversarial synthesis --- the lock identity is
fake-able by tiny rational-ratio spectra and miss-able by large-numerator organizations ---
and the competing-predictor audit \cite{c08:QG319.1} evaluates entropy, gini, exponent,
spectral gap, and lock coherence with an identical protocol on the evolving cohort, finding
the standard measures reach the same or better accuracy. This scoping limits the lock
\emph{rule} as a detector, not the lock \emph{law} as a structural result.

\section{Summary}
\label{c08:sec:locks-summary}

This chapter has established the lock law of the theory: the lock identities are Locking
projections of the spectrum, derived and not primitive (Theorem~\ref{c08:thm:locks-locking}),
universal in structure across organized domains and robust content
(Theorem~\ref{c08:thm:lock-law-universal}), with domain-specific fingerprint values
(Theorem~\ref{c08:thm:lock-values-domain}). The locks precede organization and indicate
future rigidity (Theorem~\ref{c08:thm:locks-precede}); organization emerges at a critical
transition, the locks appearing at the critical window (Theorem~\ref{c08:thm:phase-transition});
and the lock structure predicts the future organization class, though not the strength
within it (Theorem~\ref{c08:thm:blind-prediction}). The lock origin is the moment-chain
identity of the D96 spectrum (Theorem~\ref{c08:thm:lock-origin}), and the lock rule is
cohort-scoped --- a scoping that limits the detector, not the structural law. The following
chapters read the physics sectors of the theory, each as a structured read of the D96
spectrum through the operator layer and its lock structure.

\begin{thebibliography}{99}
\bibitem{c08:QG312} AT-QG Phase 312, \emph{Adversarial Spectrum Audit}, Docs/Research.
\bibitem{c08:QG313} AT-QG Phase 313, \emph{Lock Universality Audit}, Docs/Research.
\bibitem{c08:QG314} AT-QG Phase 314, \emph{Organization Predictor Audit}, Docs/Research.
\bibitem{c08:QG315} AT-QG Phase 315, \emph{Early Lock Prediction}, Docs/Research.
\bibitem{c08:QG316} AT-QG Phase 316, \emph{Organization Phase Transition Audit}, Docs/Research.
\bibitem{c08:QG318} AT-QG Phase 318, \emph{Reorganization Prediction}, Docs/Research.
\bibitem{c08:QG318.1} AT-QG Phase 318.1, \emph{Lock Origin Audit}, Docs/Research.
\bibitem{c08:QG319} AT-QG Phase 319, \emph{False Positive Audit}, Docs/Research.
\bibitem{c08:QG319.1} AT-QG Phase 319.1, \emph{Competing Predictor Audit}, Docs/Research.
\bibitem{c08:MONO007} AT-MONO007, \emph{Referee-Readiness Resolution}, Docs/Zenodo/Monograph.
\end{thebibliography}
 after the preamble
%         that defines the theorem environments (or include via the master file).
% ======================================================================

% ---- Theorem environments (define in the master preamble if not present) ----
% \newtheorem{definition}{Definition}[chapter]
% \newtheorem{lemma}[definition]{Lemma}
% \newtheorem{theorem}[definition]{Theorem}
% \newtheorem{corollary}[definition]{Corollary}
% \newtheorem{remark}[definition]{Remark}

\chapter{The Lock Law}
\label{c08:chap:locks}

\section{Introduction}
\label{c08:sec:locks-introduction}

In Chapter~\ref{c07:chap:operators} we established the operator basis of the theory: the four
operators Crowding, Compression, Beat, and Locking, derived as projections of the D96
spectrum. This chapter develops the deeper structure of the Locking operator --- the
\emph{lock law}: the lock identities of organized spectra carry a distinctive structure ---
the lock \emph{law} is universal (the presence of stable, reproducible characteristic
ratios), while the lock \emph{values} are domain-specific \cite{c08:QG313}. We establish the
temporal and structural properties of the locks --- that they precede organization
\cite{c08:QG315}, that organization emerges at a critical transition \cite{c08:QG316}, and
that the lock structure traces to the moment-chain identity \cite{c08:QG318.1}.

The lock results are established on \emph{synthetic deterministic evolving-law cohorts}
\cite{c08:QG315,c08:QG316}: the claims are cohort-scoped and the synthetic basis is disclosed,
per the referee-safe requirements of MONO007 \cite{c08:MONO007}. The honest scope of the lock
rule is reported: competing standard complexity measures match or beat the lock rule on the
evolving cohort \cite{c08:QG319.1}, and the lock rule is weak as a standalone detector under
adversarial synthesis \cite{c08:QG319}. These results scope, rather than contradict, the
lock law.

\section{Locking as an Operator}
\label{c08:sec:locking-operator}

\begin{theorem}[Locks are Locking projections]
\label{c08:thm:locks-locking}
A \emph{lock identity} is a normalized ratio of a spectrum's moments --- the moment/span,
compression/count, higher-moment, and $\sqrt{\text{moment}}/\text{span}$ ratios --- that
takes a stable, reproducible value characteristic of the spectrum's organization. The lock
identities are projections of the Locking operator: normalized ratio reads of the spectrum
\cite{c08:QG313}, consistent with the operator basis of Chapter~\ref{c07:chap:operators}
(Theorem~\ref{c07:thm:operators-derived}). They are derived reads, not primitives: by
Chapter~\ref{c07:chap:operators} (Corollary~\ref{c07:cor:operators-not-primitive}) the
operators are derived projections of the spectrum, and by Chapter~\ref{c06:chap:d96}
(Theorem~\ref{c06:thm:d96-derived}) the spectrum is the derived output of the attractor ---
the locks are thrice-removed from primitiveness.
\end{theorem}

\section{The Universal Lock Law}
\label{c08:sec:universal-law}

\begin{theorem}[The lock law is universal in structure]
\label{c08:thm:lock-law-universal}
The lock law is universal: every organized domain carries stable, reproducible lock
identities. The \emph{presence} of lock structure recurs across all organized domains.
\end{theorem}

\begin{proof}
The lock-universality audit \cite{c08:QG313} measures the normalized lock identities across
seven domains --- physics (D96 multiplicities), language (Zipf), music (harmonic
octave-class), DNA (codon inequality), software (token power law), finance (heavy tail),
networks (connectome) --- and finds stable lock structure in all seven: the lock law recurs
in every organized domain. Locks are robust content: unlike the binary operator presence,
which can be faked --- the operator basis can be partially faked by two-level crafted
spectra \cite{c08:QG312}, but the beat-identity locks are carried by none of the fakes ---
the lock identities are not trivially reproducible by adversarial spectra. The locks are
therefore the robust organization content, in contrast to the weak binary screen, and this
is the basis of the lock law's status as the deeper organization signature.
\end{proof}

\section{Domain-Specific Lock Values}
\label{c08:sec:domain-values}

\begin{theorem}[Lock values are domain-specific]
\label{c08:thm:lock-values-domain}
The lock \emph{values} are domain-specific: the characteristic ratios differ across domains,
each domain carrying its own lock values.
\end{theorem}

\begin{proof}
The lock-universality audit \cite{c08:QG313} measures the moment/span, compression/count,
higher-moment, and $\sqrt{\text{moment}}/\text{span}$ ratios per domain and finds distinct
values (at least four distinct moment/span values appear):
\begin{center}
\begin{tabular}{lcccc}
\hline
Domain & moment/span & comp./count & higher-moment & $\sqrt{\text{moment}}/\text{span}$ \\
\hline
physics & 31.7 & 2.41 & 2.96 & 21.4 \\
language & 45.0 & 180.6 & 369.8 & 5.70 \\
music & 19.7 & 22.1 & 26.3 & 4.67 \\
DNA & 36.0 & 5.67 & 6.35 & 16.3 \\
software & 9.19 & 523.2 & 877.7 & 0.84 \\
finance & 1.62 & 334.5 & 469.9 & 0.18 \\
networks & 66.3 & 3.97 & 4.02 & 33.5 \\
\hline
\end{tabular}
\end{center}
Hence the lock law has two faces: by Theorem~\ref{c08:thm:lock-law-universal} the structure
(the presence of stable ratios) is universal; by this theorem the values (the specific
ratios) are the domain's fingerprint --- universal in structure, specific in value.
\end{proof}

\section{Locks and Organization}
\label{c08:sec:locks-organization}

\begin{theorem}[Locks precede organization]
\label{c08:thm:locks-precede}
The lock identities precede mature organization: in evolving systems, the lock structure
appears before the full organization develops, and systems with early lock structure are
rigidity-committed: they reorganize less than un-locked systems under structural change.
\end{theorem}

\begin{proof}
The early-lock-prediction audit \cite{c08:QG315} evolves four systems (software history,
wiki edits, citation networks, language corpora) from flat to mature laws and measures lock
coherence and organization maturity at each stage: in most systems the lock identity reaches
half-strength at an earlier stage than the maturity does, and no system shows the locks
lagging the maturity. Hence the locks precede organization --- an early signature of the
developing structure. They are also a rigidity signal: the reorganization-prediction audit
\cite{c08:QG318} evolves systems through a reorganization and measures future topology
change --- systems with early lock coherence reorganize small (the locked group is 100\%
small-reorganization), while un-locked systems remain plastic and reorganize large.
\end{proof}

\section{The Critical Transition}
\label{c08:sec:critical-transition}

\begin{theorem}[Organization is a phase transition]
\label{c08:thm:phase-transition}
The organization structure emerges at a critical transition: the binary operator basis
completes discontinuously at a critical parameter, while the quantitative organization
grows continuously.
\end{theorem}

\begin{proof}
The organization-phase-transition audit \cite{c08:QG316} sweeps a continuous organization
parameter across the four regimes (white noise, weak, medium, strong), measuring the
operator basis, the lock coherence, and the organization maturity. The binary all-four
operator screen flips from incomplete to complete at the critical parameter $g^\ast \approx
0.31$ and persists for all stronger organization; the lock coherence emerges in the same
critical window; the maturity grows continuously. Hence the binary operator basis is a
critical order parameter --- a phase transition --- while the quantitative organization is
gradual, and the locks appear at the critical onset of organization, consistent with the
locks-precede result (Theorem~\ref{c08:thm:locks-precede}).
\end{proof}

\section{Predictive Consequences}
\label{c08:sec:predictive}

\begin{theorem}[The blind prediction protocol]
\label{c08:thm:blind-prediction}
The lock structure carries forward-looking information: the future organization class can be
predicted from the early lock structure, with the prediction rule fixed before the reveal.
\end{theorem}

\begin{proof}
The blind-organization-prediction result predicts the future high-maturity class from the
early-stage lock coherence, with the prediction rule fixed before the later stage is
revealed; the early lock presence identifies the future top-class systems exactly.
\end{proof}

The organization-predictor audit \cite{c08:QG314} computes an organization score from the
lock structure only and tests it against the eight domains: the score separates the
organized class from the unorganized class --- the organized systems lock above both
unorganized systems --- but does not rank the strength within the class: heavy-tailed
systems have large characteristic numerators that never lock onto a small fraction, scoring
below the Zipf systems despite being stronger organizations. Hence the lock values predict
the class, not the strength: the lock structure is a class-level predictor of organization,
distinguishing organized from unorganized systems early and reliably, without ranking them
internally.

\section{Validation Status}
\label{c08:sec:validation}

\begin{theorem}[The lock origin is the moment-chain identity]
\label{c08:thm:lock-origin}
The lock structure traces to the moment-chain identity of the D96 spectrum:
$\mathrm{occMom}/\Sigma m = (\Sigma m^2/\Sigma m)\cdot(\mathrm{occMom}/\Sigma m^2)$
holds exactly, linking the lock identities by an algebraic necessity.
\end{theorem}

\begin{proof}
The lock-origin audit \cite{c08:QG318.1} establishes the moment-chain (telescoping)
identity: the lock1 ratio $\mathrm{occMom}/\Sigma m = 20.0026$ equals the product of lock2
$\Sigma m^2/\Sigma m = 2.4105$ and lock3 $\mathrm{occMom}/\Sigma m^2 = 8.2980$, exactly.
The locks are therefore ratios of one moment hierarchy, linked by an algebraic necessity ---
the common source of lock formation: a resonance fixed point in structure, D96-specific in
value.
\end{proof}

The lock rule is cohort-scoped: the false-positive audit \cite{c08:QG319} finds the lock
rule weak as a standalone detector under adversarial synthesis --- the lock identity is
fake-able by tiny rational-ratio spectra and miss-able by large-numerator organizations ---
and the competing-predictor audit \cite{c08:QG319.1} evaluates entropy, gini, exponent,
spectral gap, and lock coherence with an identical protocol on the evolving cohort, finding
the standard measures reach the same or better accuracy. This scoping limits the lock
\emph{rule} as a detector, not the lock \emph{law} as a structural result.

\section{Summary}
\label{c08:sec:locks-summary}

This chapter has established the lock law of the theory: the lock identities are Locking
projections of the spectrum, derived and not primitive (Theorem~\ref{c08:thm:locks-locking}),
universal in structure across organized domains and robust content
(Theorem~\ref{c08:thm:lock-law-universal}), with domain-specific fingerprint values
(Theorem~\ref{c08:thm:lock-values-domain}). The locks precede organization and indicate
future rigidity (Theorem~\ref{c08:thm:locks-precede}); organization emerges at a critical
transition, the locks appearing at the critical window (Theorem~\ref{c08:thm:phase-transition});
and the lock structure predicts the future organization class, though not the strength
within it (Theorem~\ref{c08:thm:blind-prediction}). The lock origin is the moment-chain
identity of the D96 spectrum (Theorem~\ref{c08:thm:lock-origin}), and the lock rule is
cohort-scoped --- a scoping that limits the detector, not the structural law. The following
chapters read the physics sectors of the theory, each as a structured read of the D96
spectrum through the operator layer and its lock structure.

\begin{thebibliography}{99}
\bibitem{c08:QG312} AT-QG Phase 312, \emph{Adversarial Spectrum Audit}, Docs/Research.
\bibitem{c08:QG313} AT-QG Phase 313, \emph{Lock Universality Audit}, Docs/Research.
\bibitem{c08:QG314} AT-QG Phase 314, \emph{Organization Predictor Audit}, Docs/Research.
\bibitem{c08:QG315} AT-QG Phase 315, \emph{Early Lock Prediction}, Docs/Research.
\bibitem{c08:QG316} AT-QG Phase 316, \emph{Organization Phase Transition Audit}, Docs/Research.
\bibitem{c08:QG318} AT-QG Phase 318, \emph{Reorganization Prediction}, Docs/Research.
\bibitem{c08:QG318.1} AT-QG Phase 318.1, \emph{Lock Origin Audit}, Docs/Research.
\bibitem{c08:QG319} AT-QG Phase 319, \emph{False Positive Audit}, Docs/Research.
\bibitem{c08:QG319.1} AT-QG Phase 319.1, \emph{Competing Predictor Audit}, Docs/Research.
\bibitem{c08:MONO007} AT-MONO007, \emph{Referee-Readiness Resolution}, Docs/Zenodo/Monograph.
\end{thebibliography}
